%% ma: L12-B08-0088
%% lop: 12
%% chuong: 2
%% bai: 8
%% dang: 12.8.42
%% loai: TN4
%% muc: TH
%% dap_an: "A"
%% nguon: "CĐ2.2.pdf giáo viên gửi 27/10/2026 (M07, Lê Minh Khôi – Nguyễn Anh Khôi, Chương II Chủ đề 2 Tọa độ của vectơ trong không gian), Câu 22"
%% kien_thuc: "Bài 7 (hai vectơ bằng nhau khi các tọa độ tương ứng bằng nhau)"
%% trang_thai: chinh-thuc
%% ghi_chu: "Giáo viên đã duyệt 27/10/2026. Đã sửa giả thiết m, n thuộc Z thành m, n thuộc R (m = 1/3 không nguyên). Đề gốc ghi m, n thuộc Z nhưng m = 1/3 không nguyên; đã sửa thành m, n thuộc R (đáp án không đổi)"
\begin{ex}
Trong không gian $Oxyz$, cho hai vectơ $\vec u=(3;2;5)$ và $\vec v=(3m+2;2;6-n)$ với $m,n\in\mathbb R$. Biết rằng $\vec u=\vec v$. Tính giá trị của biểu thức $Z=2025m+2027n$.
\choice{\True $Z=2702$}{$Z=2072$}{$Z=4052$}{$Z=4050$}
\loigiai{$\vec u=\vec v\Leftrightarrow3m+2=3$ và $6-n=5\Leftrightarrow m=\dfrac13$, $n=1$.

$Z=675+2027=2702$.}
\end{ex}
